% DOC NO. RL-Z0-IDEAS · REV. A
%
% This file IS the document. Edit it directly - ripdoc compiles it
% as it stands and stamps the provenance line at build time.
%
%   ripdoc build --only docs-ideas
%
% Promoted from docs/ideas.md on 2026-09-07 by `ripdoc promote`.
\documentclass[fontpath=fonts/,logopath=logo/]{riposte-doc}
\usepackage{riposte-md}

\title{CHANNEL Z0 — Ideas \& Roadmap (Station Bulletin № 2)}
\author{Riposte Laboratories}
\date{}
\ripdocno{RL-Z0-IDEAS}
\riprev{A}
\ripstrapline{\NoCaseChange{\riphi{A local channel is a premise, not a product; it can keep growing without ever getting bigger. This is the writers' room — what's already on the air, and what's sketched on the whiteboard. Steal freely; the plumbing is MIT.}}}
\riprunningtitle{CHANNEL Z0 — Ideas \& Roadmap (Station Bulletin № 2)}

\begin{document}

\ripmakehead


\ripsection{\ripmdhead{Shipped in this pass}}

These went from idea to implementation and are in the repo now.

\begin{ripbullets}
\item \textbf{The live marquee.} The storefront's "NOW SHOWING" used to read from a hand-typed schedule. Now \ripmdpath{\ripmono{playout/\allowbreak{}nowplaying.\allowbreak{}py}} reads what ErsatzTV is \riphi{actually} airing from its XMLTV guide and writes it into Owncast's stream title; the site shows the real program, live. Ships as an optional service in \ripmdpath{\ripmono{playout/\allowbreak{}compose.\allowbreak{}yml}} (\ripmono{-\allowbreak{}-\allowbreak{}profile nowplaying}). This is the difference between a webpage about a channel and a webpage \riphi{tuned to} one.
\item \textbf{The whole home side in one file.} \ripmdpath{\ripmono{playout/\allowbreak{}compose.\allowbreak{}yml}} containerizes master control — ErsatzTV, the uplink relay (an \ripmono{ffmpeg -\allowbreak{}c copy} loop that replaces the systemd unit), and the marquee — so it drops onto a Proxmox LXC or TrueNAS box without ceremony. See \ripdocref{\ripmono{docs/\allowbreak{}self-\allowbreak{}hosting.\allowbreak{}md}}{RL-Z0-SELF-HOSTING}.
\item \textbf{Station identity, generated.} The build guide asked for bumpers, a channel bug, and a technical-difficulties slate, then never made them. Now:
\begin{ripbullets}
\item \ripmdpath{\ripmono{tools/\allowbreak{}make-\allowbreak{}ident.\allowbreak{}sh}} — the "NOW WATCHING CHANNEL Z0" bumpers the Station ID filler rotates.
\item \ripmdpath{\ripmono{tools/\allowbreak{}make-\allowbreak{}slate.\allowbreak{}sh}} — full-screen slate cards (technical difficulties, sign-off, please-stand-by), 1080p, on-brand.
\item \ripmdpath{\ripmono{tools/\allowbreak{}make-\allowbreak{}bug.\allowbreak{}sh}} — the transparent channel bug PNG for ErsatzTV's watermark. All three speak the storefront's language — monospace on ink, one red, the RL-Z0 designation — so the on-screen channel and the website look like one station.
\end{ripbullets}

\item \textbf{The bookends of the broadcast day.} The two pieces of furniture that top and tail a day of television, generated on brand:
\begin{ripbullets}
\item \ripmdpath{\ripmono{tools/\allowbreak{}make-\allowbreak{}testcard.\allowbreak{}sh}} — the station test card. Not raw \ripmono{smptehdbars} but the card \riphi{with the station's name on it}: geometry grid, colour bars, a greyscale step wedge, corner castellations, a centre identity panel, and a 1 kHz line-up tone. The pre-06:00 sign-on fill, and a picture the tower can align to.
\item \ripmdpath{\ripmono{tools/\allowbreak{}make-\allowbreak{}signoff.\allowbreak{}sh}} — the nightly sign-off. The "THIS CONCLUDES OUR BROADCAST DAY / GOODNIGHT, LOCALS" card dissolving into a tail of colour bars, one file for the 00:00 SIGN-OFF item. The seed of the sign-off anthem below, close at hand.
\end{ripbullets}

\item \textbf{Pre-flight for spots.} \ripmdpath{\ripmono{tools/\allowbreak{}check-\allowbreak{}ad.\allowbreak{}sh}} reads a submission and tells you its length, codecs, and true loudness before you accept it — the screening step that pairs with \ripmono{normalize-\allowbreak{}ad.\allowbreak{}sh}.
\item \textbf{The broadcast week, written down.} The build guide gave master control a one-paragraph "starting rhythm" and an empty ErsatzTV; now \ripdocref{\ripmono{docs/\allowbreak{}programming.\allowbreak{}md}}{RL-Z0-PROGRAMMING} is the whole week — five dayparts, a protected 19:00 flagship, themed nights (Atomic Tuesday, Friday Night Feature, the Saturday double bill), the show bible for what's ours, and the exact ErsatzTV mapping (collections, fixed-start anchors, the weekend schedule). The storefront's program grid (\ripmdpath{\ripmono{site/\allowbreak{}index.\allowbreak{}html}}) is now its machine-readable twin: a real seven-day \ripmono{WEEK} with day tabs, so viewers can flip through the week, not just today. Guide and guide-on-the-wall, in sync.
\item \textbf{On Cloudflare.} The storefront ships to Cloudflare (Workers Static Assets) (\ripmono{ch0.\allowbreak{}ripostelabs.\allowbreak{}xyz} today, \ripmono{channelz0.tv} when purchased) with tidy short links (\ripmono{/watch}, \ripmono{/lab}) and security headers. The tower stays on the VPS.
\end{ripbullets}


\ripsection{\ripmdhead{On the workbench}}

Sketches worth building next, roughly in order of bang-for-effort.


\ripsubsection{\ripmdhead{Programming}}

\begin{ripbullets}
\item \textbf{The Neighbourhood Desk, for real.} A \ripmono{bulletin.json} the storefront renders as a scrolling ticker — lost cats, garage sales, the school play. Locals are the content; the ticker is the cheapest possible community TV.
\item \textbf{Ask the alien.} A submission form (a Cloudflare Worker → email or R2) where locals leave questions for the Ground Zero correspondent to ask the neighbourhood via baguette. Turns viewers into a writers' room.
\item \textbf{The sign-off anthem.} A 60–90s nightly close — a slow pan over the lab, the RL-Z0 card, a bit of music you own — generated once and aired at 00:00 before the colour bars. \ripmono{make-slate.sh} is the seed; this is its cinematic cousin.
\item \textbf{Emergency Broadcast crawl.} A parody EAS: the checker band, the two-tone, a crawl of gloriously low-stakes alerts ("A CASSEROLE HAS BEEN LEFT UNATTENDED"). Aired rarely, on purpose.
\end{ripbullets}


\ripsubsection{\ripmdhead{On-air polish}}

\begin{ripbullets}
\item \textbf{Lower-thirds kit.} An ffmpeg/overlay template for Ground Zero — name, location, "DEFINITELY HUMAN" — so field interviews get a broadcast chyron.
\item \textbf{The QR corner.} A tiny static QR to the storefront, burned into the bug area during idents, so anyone who wanders past a TV can find the stream.
\end{ripbullets}


\ripsubsection{\ripmdhead{The site}}

\begin{ripbullets}
\item \textbf{Full week guide.} ErsatzTV already serves XMLTV; a small build step could turn a week of it into a printable TV-guide page (very on-brand as a zine).
\item \textbf{Receiver map / "who's tuned in."} Owncast reports viewer counts; a sparkline of the day's audience makes the "someone's always watching" promise visible.
\end{ripbullets}


\ripsubsection{\ripmdhead{The bigger swing}}

\begin{ripbullets}
\item \textbf{Channel Z1.} The architecture is per-channel, not per-station. A second ErsatzTV channel + a second Owncast is a sibling network (Z1: overnight ambient? all-cartoons?). The storefront becomes a dial.
\end{ripbullets}

\ripdashrule

\riphi{Ideas are cheap; airtime is free while we're small. Build the ones that make the neighbourhood lean in.}

\ripend

\end{document}
